\documentclass[11pt,a4paper]{article}
\usepackage[utf8]{inputenc}
\usepackage[vietnamese]{babel}
\usepackage{amsmath,amssymb}
\usepackage{geometry}
\geometry{a4paper, total={180mm,260mm}, left=15mm, right=15mm, top=18mm, bottom=18mm}

\usepackage{tikz}
\usetikzlibrary{shapes.geometric, arrows.meta, positioning, calc}
\tikzset{
  arr/.style={-{Stealth[scale=0.8]}},
  every node/.append style={align=center}
}
\usepackage{tcolorbox}
\tcbuselibrary{skins,breakable}
\usepackage{fontawesome}
\usepackage{fancyhdr}
\usepackage{enumitem}

\pagestyle{fancy}
\fancyhf{}
\renewcommand{\headrulewidth}{0.5pt}
\renewcommand{\footrulewidth}{0.5pt}
\fancyhead[L]{\small\bfseries\color{violet!85!black}\faBook\ SỔ TAY KIẾN THỨC TOÁN 2}
\fancyhead[R]{\small\bfseries\color{violet!85!black}Lớp Toán Thầy Lĩnh -- loptoanthaylinh.vn}
\fancyfoot[C]{\small\color{gray}-- Sổ tay kiến thức: Bài 13. Bài toán về nhiều hơn, ít hơn một số đơn vị --}

\setlength{\parindent}{0pt}
\setlength{\parskip}{0pt}

\begin{document}

% --- BANNER BÀI HỌC ---
\begin{tcolorbox}[
  enhanced,
  colback=violet!8!white,
  colframe=violet!70!black,
  arc=3mm, boxrule=1.2pt,
  left=3mm, right=3mm, top=2mm, bottom=2mm,
  halign=center
]
  {\LARGE\bfseries\color{violet!90!black}\faBookmark\ BÀI 13: BÀI TOÁN VỀ NHIỀU HƠN, ÍT HƠN MỘT SỐ ĐƠN VỊ}\\[1mm]
  {\itshape\footnotesize\color{black!75} Phương pháp tìm số lớn khi biết số bé và phần hơn; tìm số bé khi biết số lớn và phần kém trên 1 trang A4}
\end{tcolorbox}

\vspace{-2mm}

\noindent
\begin{minipage}[t]{0.49\textwidth}\vspace{0pt}
\begin{tcolorbox}[
  enhanced,
  colback=violet!4!white,
  colframe=violet!75!black,
  arc=2mm, boxrule=0.8pt,
  left=2mm, right=2mm, top=1.8mm, bottom=1.8mm,
  title={\footnotesize\bfseries\color{white}\faKey\ 1. Dạng 1: Bài toán về nhiều hơn}
]
{\footnotesize
\begin{itemize}\setlength{\itemsep}{1pt}\setlength{\leftskip}{-2.5mm}
  \item \textbf{Đặc điểm:} Đề bài cho biết \textbf{số bé} và \textbf{phần nhiều hơn}, yêu cầu tìm \textbf{số lớn}.
  \item \textbf{Quy tắc tính:}\\
  Muốn tìm số lớn, ta lấy số bé \textbf{cộng} với phần nhiều hơn:
  $$\mathbf{\text{Số lớn} = \text{Số bé} + \text{Phần nhiều hơn}}$$
  \item \textbf{Ví dụ:} An có $8$ viên bi, Bình có nhiều hơn An $5$ viên bi. Hỏi Bình có bao nhiêu viên bi?\\
  $\implies$ Bình có: $8 + 5 = \mathbf{13}$ (viên bi).
\end{itemize}
}
\end{tcolorbox}
\end{minipage}
\hfill
\begin{minipage}[t]{0.49\textwidth}\vspace{0pt}
\begin{tcolorbox}[
  enhanced,
  colback=teal!4!white,
  colframe=teal!75!black,
  arc=2mm, boxrule=0.8pt,
  left=2mm, right=2mm, top=1.8mm, bottom=1.8mm,
  title={\footnotesize\bfseries\color{white}\faSliders\ 2. Dạng 2: Bài toán về ít hơn}
]
{\footnotesize
\begin{itemize}\setlength{\itemsep}{1pt}\setlength{\leftskip}{-2.5mm}
  \item \textbf{Đặc điểm:} Đề bài cho biết \textbf{số lớn} và \textbf{phần ít hơn}, yêu cầu tìm \textbf{số bé}.
  \item \textbf{Quy tắc tính:}\\
  Muốn tìm số bé, ta lấy số lớn \textbf{trừ} đi phần ít hơn:
  $$\mathbf{\text{Số bé} = \text{Số lớn} - \text{Phần ít hơn}}$$
  \item \textbf{Ví dụ:} Cành trên có $14$ quả cam, cành dưới có ít hơn cành trên $6$ quả. Hỏi cành dưới có bao nhiêu quả?\\
  $\implies$ Cành dưới có: $14 - 6 = \mathbf{8}$ (quả).
\end{itemize}
}
\end{tcolorbox}
\end{minipage}

\vspace{-1.5mm}
\begin{tcolorbox}[
  enhanced,
  colback=orange!4!white,
  colframe=orange!80!black,
  arc=2mm, boxrule=0.8pt,
  left=2mm, right=2mm, top=1.5mm, bottom=1.5mm,
  title={\footnotesize\bfseries\color{white}\faSitemap\ 3. Sơ đồ đoạn thẳng so sánh: Tìm Số Lớn và Tìm Số Bé}
]
\centering
\begin{tikzpicture}[x=1cm, y=0.8cm]
  % Dạng nhiều hơn
  \node[font=\scriptsize\bfseries, violet!80!black] at (2.5, 1.6) {TÌM SỐ LỚN: Phép tính CỘNG (+)};
  \draw[thick, teal!70!black, fill=teal!15] (0, 0.9) rectangle (2.5, 1.3) node[midway, font=\tiny\bfseries] {Số bé};
  \draw[thick, violet!70!black, fill=violet!15] (0, 0.2) rectangle (4, 0.6);
  \draw[dashed, red!80!black] (2.5, 0.9) -- (2.5, 0.2);
  \node[font=\tiny\bfseries, red!85!black] at (3.25, 0.4) {Nhiều hơn};
  \node[font=\tiny\bfseries] at (1.25, 0.4) {Bằng số bé};

  % Dạng ít hơn
  \node[font=\scriptsize\bfseries, teal!80!black] at (8, 1.6) {TÌM SỐ BÉ: Phép tính TRỪ (-)};
  \draw[thick, violet!70!black, fill=violet!15] (5.5, 0.9) rectangle (10, 1.3) node[midway, font=\tiny\bfseries] {Số lớn};
  \draw[thick, teal!70!black, fill=teal!15] (5.5, 0.2) rectangle (8, 0.6) node[midway, font=\tiny\bfseries] {Số bé};
  \draw[dashed, red!80!black] (8, 0.9) -- (8, 0.2);
  \node[font=\tiny\bfseries, red!85!black] at (9, 1.1) {Ít hơn};
\end{tikzpicture}
\end{tcolorbox}

\vspace{-1.5mm}
\noindent
\begin{minipage}[t]{0.49\textwidth}\vspace{0pt}
\begin{tcolorbox}[
  enhanced,
  colback=green!4!white,
  colframe=green!60!black,
  arc=2mm, boxrule=0.8pt,
  left=2mm, right=2mm, top=1.8mm, bottom=1.8mm,
  title={\footnotesize\bfseries\color{white}\faCheckCircle\ 4. Ví dụ thực hành chuẩn mực}
]
{\footnotesize
\textbf{Bài toán mẫu (Dạng ít hơn):}\\
Anh có $15$ cái kẹo, em có ít hơn anh $6$ cái kẹo. Hỏi em có bao nhiêu cái kẹo?\\
\textit{Bài giải:}\\
Em có số cái kẹo là:
$$15 - 6 = 9\text{ (cái kẹo)}$$
Đáp số: $9$ cái kẹo.
}
\end{tcolorbox}
\end{minipage}
\hfill
\begin{minipage}[t]{0.49\textwidth}\vspace{0pt}
\begin{tcolorbox}[
  enhanced,
  colback=red!4!white,
  colframe=red!75!black,
  arc=2mm, boxrule=0.8pt,
  left=2mm, right=2mm, top=1.8mm, bottom=1.8mm,
  title={\footnotesize\bfseries\color{white}\faExclamationTriangle\ 5. Mẹo tránh bẫy trắc nghiệm}
]
{\footnotesize
\begin{itemize}\setlength{\itemsep}{1pt}\setlength{\leftskip}{-2.5mm}
  \item \textbf{Phân biệt hai dạng câu hỏi:}\\
  - \textit{"Hỏi hơn/kém bao nhiêu?"} $\to$ Tìm phần chênh lệch $\to$ \textbf{Làm phép TRỪ}.\\
  - \textit{"Hỏi đối tượng thứ hai có bao nhiêu?"} $\to$ Nhìn từ khóa "nhiều hơn" ($\to$ CỘNG), "ít hơn" ($\to$ TRỪ).
  \item Luôn vẽ sơ đồ đoạn thẳng ra nháp để không nhầm chiều so sánh.
\end{itemize}
}
\end{tcolorbox}
\end{minipage}


\end{document}
